\documentclass[a4paper]{article}
\usepackage{amsfonts}
\usepackage{amsmath}
\usepackage{amssymb}
\usepackage{graphicx}     

\newcommand{\dint}{\displaystyle\int}
\newcommand{\llim}{\lim\limits}

\begin{document}
  
\noindent \large Auteur: Oubayy Ahale          \\
\noindent \large Studierichting: Informatica   \\
\noindent \large Oefeningenreeks 2B nr. 6.4    \\

\medskip

\normalsize

Bewijs dat: $ \operatorname{Bgtan} \dfrac{1}{2} + \operatorname{Bgtan} \dfrac{1}{5} + \operatorname{Bgtan} \dfrac{1}{8} = \dfrac{\pi}{4} $\\

\medskip

Stel $ \Sigma = \operatorname{Bgtan} \dfrac{1}{2} $, $ \alpha = \operatorname{Bgtan} \dfrac{1}{5} $, $ \beta = \operatorname{Bgtan} \dfrac{1}{8} $\\

$ \Rightarrow \tan(\Sigma) = \dfrac{1}{2}, \tan(\alpha) = \dfrac{1}{5}, \tan(\beta) = \dfrac{1}{8} $\\

\medskip

$ \Rightarrow \Sigma + \alpha + \beta = \dfrac{\pi}{4} $\\

$ \Rightarrow \tan(\Sigma + (\alpha + \beta)) = \tan\left(\dfrac{\pi}{4}\right) = 1 $\\

Linkerzijde $ = \dfrac{\tan \Sigma + \tan (\alpha + \beta)}{1 - \tan \Sigma \cdot \tan (\alpha + \beta)} $\\

$ = \dfrac{\tan \Sigma + \dfrac{\tan \alpha + \tan \beta}{1 - \tan \alpha \cdot \tan \beta}}{1 - \tan \Sigma \cdot \dfrac{\tan \alpha + \tan \beta}{1 - \tan \alpha \cdot \tan \beta}} $\\

\medskip

$ = \dfrac{\dfrac{1}{2} + \dfrac{\dfrac{1}{5} + \dfrac{1}{8}}{1 - \dfrac{1}{5} \cdot \dfrac{1}{8}} }  {1 - \dfrac{1}{2} \cdot \dfrac{\dfrac{1}{5} + \dfrac{1}{8}}{1 - \dfrac{1}{5} \cdot \dfrac{1}{8}}} $\\

\medskip

$ = \dfrac{\dfrac{1}{2} + \dfrac{1}{3}}{1 - \dfrac{1}{2} \cdot \dfrac{1}{3}} $\\

\medskip

$ = \dfrac{\dfrac{3}{6} + \dfrac{2}{6}}{\dfrac{6}{6} - \dfrac{1}{6}} $\\

\medskip

$ = \dfrac{\dfrac{5}{6}}{\dfrac{5}{6}} = 1 $\\

\medskip

Rechterzijde $ = \tan \dfrac{\pi}{4} = 1 $\\

$ \Rightarrow $ L.L $=$ R.L \\

Q.E.D.

\begin{thebibliography}{999}
%\bibitem{a} Referentie
\end{thebibliography}
\end{document} 